%\ifodd\value{page}\hbox{}\newpage\fi
\chapter*{Śrī Lalitā Sahasranāma — Śloka 12}
\addcontentsline{toc}{chapter}{Śloka 12 - Nomi 26,27}

\begin{sloka}
{\sanskritfont
\begin{tabular}{c}
अनाकलितसादृश्यचिबुकश्रीविराजिता ।\\
कामेशबद्धमाङ्गल्यसूत्रशोभितकन्धरा ॥ १२ ॥
\end{tabular}

\vspace{1.5em}

{\middlesize \iastfont
anākalita-sādṛśya-cibuka-śrī-virājitā |\\
kāmeśa-baddha-māṅgalya-sūtra-śobhita-kandharā ||12||}

\vspace{1em}

{\middlesize
\anvaya{%
अनाकलितसादृश्यचिबुकश्रीविराजिता ।
कामेशबद्धमाङ्गल्यसूत्रशोभितकन्धरा ।
}
}

\middlesize \iastfont \itmean{%
«Il cui mento risplende di una bellezza senza paragone;  
il cui collo è adornato dal filo augurale annodato da \textit{Kāmeśa}.»%
}

}
\end{sloka}

\wordmeaning{%
\wordline{अनाकलित-सादृश्य-चिबुक-श्री-विराजिता}{anākalita-sādṛśya-cibuka-śrī-\\ virājitā}
  { il cui mento risplende di una bellezza incomparabile, priva di similitudine;}%
\wordline{कामेश-बद्ध-माङ्गल्य-सूत्र-शोभित-कन्धरा}{kāmeśa-baddha-māṅgalya-sūtra-\\ śobhita-kandharā}
  { il cui collo è reso splendente dal filo augurale (\textit{māṅgalya-sūtra}), il filo nuziale tradizionale, annodato da Kāmeśa.}%
}

Lo \textit{Śloka} 8 prosegue l’itinerario discendente della contemplazione soffermandosi sul mento e sul collo, punti di raccordo tra volto, parola e corpo. Il mento (\textit{cibuka}) è qualificato come portatore di una bellezza \textit{anākalita-sādṛśya}, cioè priva di confronto possibile: non una forma che rinvia ad altro, ma una presenza che si impone per immediatezza e compiutezza.

Il collo (\textit{kandharā}), ornato dal filo augurale (\textit{māṅgalya-sūtra}) annodato da \textit{Kāmeśa}, introduce una dimensione relazionale e rituale. Non si tratta di un ornamento estetico, ma del segno di un vincolo cosmico: l’unione tra coscienza (\textit{Śiva}), qui nella forma di \textit{Kāmeśa}, e potenza (\textit{Śakti}) resa visibile come ordine, protezione e fecondità. Il collo diviene così luogo di passaggio tra interiorità e manifestazione, tra parola non detta e parola pronunciata.

Dal punto di vista simbolico, il mento stabilisce la chiusura del volto come unità percettiva, mentre il collo ne consente l’apertura verso il corpo e l’azione. La bellezza qui non è dispersione né compiacimento, ma equilibrio strutturale: ciò che tiene insieme senza rigidità, ciò che sostiene senza costringere.

Questo \textit{śloka} contiene due \textit{nāma} propriamente detti. \textbf{\textit{Anākalita-sādṛśya-cibuka-śrī-virājitā}} afferma la bellezza del mento come qualità non comparativa, segno di una forma che non deriva da modelli ma si impone come norma di se stessa. \textbf{\textit{Kāmeśa-baddha-māṅgalya-sūtra-śobhita-kandharā}} descrive il collo come sede del vincolo augurale, indicando che la bellezza di \textit{Lalitā} è inseparabile dalla relazione cosmica che la fonda.

In questo \textit{śloka}, la bellezza si manifesta come continuità e legame: non frammento isolato, ma forma che unisce, connette e sostiene. È la grazia di ciò che tiene insieme il visibile e l’invisibile, il volto e il corpo, la coscienza e il mondo.
